\subsection{A prospective token-budget sensitivity}\label{sec:budget}
Before the reachable-state proof, we prospectively specified
that reducing only the batch budget would make chunked-FCFS and decode-first
execute different per-request spans in at least one matched condition.
The predicted mechanism was prefill work exhausting the budget before
later resident decodes. The unmodified engine requires every whole prompt
to fit the budget for all policies. Full-corpus feasibility requires at least
29,067 tokens, above the 4,096-token chunked-work bound. No intact Azure or
Bailian window fits below 4,096 tokens; using them would require another
contract or input change.
We therefore selected the lexicographically first three frozen BurstGPT
windows with maximum prompts strictly above 512 and at most 1,024 tokens:
\texttt{pilot\_v2\_w16}, \texttt{w23}, and \texttt{w28}, with maxima
1,005, 790, and 648. This input-selected stratum permits multi-batch
prefills but is not representative of the full corpus.

All 200 requests, frozen seeds and arrival/deadline transformations remain
unchanged. We retain LOW and HIGH\_PRESSURE, eight resident slots, and
all four adaptations. Budgets are the original 32,768, then 2,048 and
1,024 tokens, corresponding to four and two 512-token chunks versus the
4,096-token potential work bound. Timing, reservations, horizon and policy
code are fixed. A 12-cell feasibility pilot precedes the remaining 60 cells;
all 72 pass the independent event referee, complete all 14,400 request
executions, and have independently recomputed terminal metrics. The 24
baseline cells exactly reproduce their archived metrics and terminal states.

The primary hypothesis is rejected for this design: normalized token-span
journals and terminal states agree for chunked-FCFS and decode-first in
all 12 perturbed matches, and three distinct terminal behaviors remain in
each of the 18 budget--condition combinations. Both policies serve every
ready resident decode in these runs; budget exhaustion changes prefill
progress without separating their transitions. Theorem~\ref{thm:equivalence} now explains this agreement: a binding
budget can truncate only the youngest resident prefill, leaving every
ready decode served. The empirical checks alone would not prove this claim. Span-list ordering differs in four of six matches at every budget,
but causes no state or timing difference.

The perturbation does affect execution. Across all four labels, exactly
full-budget batches increase from 0/25,983 to 3/25,986 and 36/25,984 as the
budget decreases. These counts include the equivalent labels and exclude
whole-prompt packing that leaves unusable residual budget. At 1,024 tokens,
SLO-gated's prefill-before-residual-decode rule omits 42 ready-decode
opportunities out of 35,976; its maximum anchor age reaches 2\,ms, still below the
unchanged 495\,ms urgency threshold. The gate never activates.
Equation~\ref{eq:urgency-bound} strengthens this observation: total prompt
work of 39,300, 52,250 and 44,765 tokens gives respective 38, 51 and
44\,ms upper bounds at budget 1,024, ruling out first activation in
each window. Independent reconstruction of all 18 saved SLO-gated cells
finds zero urgent decisions in 107,844 ready-decode opportunities.
Table~\ref{tab:budget} shows modest, sometimes nonmonotone terminal effects.
Whole-FCFS changes which whole prompts fit and hence which batches defer
decodes; a smaller budget need not worsen every completion time.

\begin{table}[!htbp]
\centering\small\setlength{\tabcolsep}{3pt}
\caption{Budget sensitivity: minimum and maximum paired changes from
32,768 tokens across six fixed window--setting conditions. Latency columns
are milliseconds; throughput is model tokens/s. Decode-first equals
chunked-FCFS in every row. Ranges describe these conditions, not confidence
intervals. All other contract parameters remain fixed.}
\label{tab:budget}\begin{tabular}{@{}rlrrr@{}}\toprule
Budget & Adaptation & $\Delta$mean & $\Delta$p95 & $\Delta$throughput \\ \midrule
2048 & Whole-FCFS & [0.000, 0.565] & [0.000, 1.000] & [-10.931, 0.000] \\
2048 & Chunked-FCFS & [0.000, 0.135] & [0.000, 0.000] & [0.000, 0.000] \\
2048 & SLO-gated & [0.000, 0.015] & [0.000, 0.000] & [0.000, 0.000] \\
1024 & Whole-FCFS & [-2.245, 2.560] & [-7.000, 4.000] & [-18.604, 18.332] \\
1024 & Chunked-FCFS & [0.000, 0.520] & [0.000, 1.897] & [0.000, 0.000] \\
1024 & SLO-gated & [0.000, 0.365] & [0.000, 2.000] & [-4.503, 0.000] \\
\bottomrule\end{tabular}

\end{table}

No strict-10\% policy-pair preference reverses across budgets
($0/576$ eligible comparisons: six conditions, two perturbations, six
policy pairs, eight endpoints). With the existing numerical-tie tolerance,
5/96 top sets change, all involving TPOT as SLO-gated leaves a former tie;
its largest mean-TPOT increase is only 5.87\,$\mu$s. Whole-FCFS deadline
attainment changes by $-0.005$ to $+0.010$, and SLO-gated by $-0.005$ to
zero. All throughput changes are below 0.4\% in magnitude. These bounded
findings preserve the tested equivalence while exposing small execution
changes; they do not retest cross-source means. Equivalence outside
this stratum follows from Theorem~\ref{thm:equivalence} only within its assumptions. Three selected, temporally related windows support
paired descriptive ranges, not population confidence intervals. The original
C01 and reduced-comparison results remain separate and unchanged.
